\documentclass{article}
\title{Tiny Paper: A Minimal Example}
\author{QiFa Fixtures}
\date{2026}

\begin{document}
\maketitle

\section{Introduction}
This tiny paper exists only as a fixture for QiFa tests. It proposes a
scoring function that ranks candidate answers by a simple linear rule.

\section{Method}
Given a candidate $x$ with features $f(x)$, the score is
\begin{equation}
s(x) = w^\top f(x) + b .
\label{eq:score}
\end{equation}
Figure~\ref{fig:arch} shows the two-stage pipeline: features are extracted
first, then scored. Table~\ref{tab:results} reports the fixture numbers.

\begin{figure}[h]
  \centering
  \fbox{features -> score -> rank}
  \caption{Two-stage scoring pipeline.}
  \label{fig:arch}
\end{figure}

\begin{table}[h]
  \centering
  \begin{tabular}{lc}
    Method & Score \\
    Baseline & 0.50 \\
    Ours & 0.62 \\
  \end{tabular}
  \caption{Fixture results.}
  \label{tab:results}
\end{table}

\end{document}
